\section{Dark Mode Implementation --- \texorpdfstring{\lstinline|next-common|}{next-common}}
\label{sec:dark-mode}

\subsection{Context}

\lstinline|next-common| is the shared front-end foundation used across every app-fo at Rakuten France. It ships common React components, hooks, design tokens and tooling that downstream apps consume as versioned npm packages.

This ticket is an \textbf{introductory task}: it forces the assignee to touch every corner of the monorepo --- theme, tokens, components, SSR, Storybook, build pipeline --- and to interact with every downstream consumer.

\textbf{Motivation:}
\begin{itemize}
    \item Many users already run their OS/browser in dark mode by default.
    \item Younger cohorts (the KRS growth target) actively prefer dark UI.
    \item It brings the platform in line with peer e-commerce sites and modern UX standards.
\end{itemize}

\subsection{Phase 01 --- Research \& Technical Planning}

\subsubsection{What is dark mode / why does it matter}

\begin{itemize}
    \item OS-level and browser-level color scheme preferences that flip UI palettes from light to dark.
    \item Detectable in CSS via \lstinline|prefers-color-scheme: dark| and in JS via \lstinline|window.matchMedia('(prefers-color-scheme: dark)')|.
    \item E-commerce specifics: product images shot on white backgrounds, price/CTA contrast, brand red on dark surfaces, seller trust signals (badges, ratings) must remain legible.
\end{itemize}

\subsubsection{Technical decisions}

\begin{table}[h]
    \centering
    \caption{Technical decisions for the dark mode implementation}
    \label{tab:dark-mode-decisions}
    \begin{tabular}{@{}p{2.4cm}p{3.6cm}p{4.4cm}p{3.6cm}@{}}
        \toprule
        \textbf{Concern} & \textbf{Options considered} & \textbf{Choice on this branch} & \textbf{Where it lives} \\
        \midrule
        State management & React Context, Zustand, Redux & \textbf{React Context} (\lstinline|AppThemeContext|) --- no need for a store, theme is a single boolean-ish value & \lstinline|packages/mui/src/components/| \lstinline|ThemeContextProvider/themeContext.ts| \\
        CSS strategy & \lstinline|.dark-mode| class, MUI-only, \lstinline|:root[data-theme]| & \textbf{\lstinline|:root[data-theme='dark'|dark|'light']|} --- one selector drives both raw CSS and MUI & \lstinline|packages/global-css/src/styles/palettes/*.css| \\
        Persistence & sessionStorage, cookie, localStorage & \textbf{\lstinline|localStorage|} key \lstinline|"theme"| --- survives reloads, no server round-trip & \lstinline|ThemeContextProvider| + \lstinline|ThemeInitScript| \\
        MUI $\leftrightarrow$ CSS-vars sync & Duplicate palette in JS, MUI v7 CssVarsProvider, JS tokens $\to$ CSS vars & \textbf{JS tokens re-declared as \lstinline|var(--*)| references + \lstinline|cssVariables: true| in \lstinline|createTheme|} so both worlds resolve to the same CSS variables & \lstinline|packages/global-css/src/colorPalette.ts| + \lstinline|packages/mui/src/theme.tsx| \\
        FOUC (Flash of Unstyled Content) & CSS-only, cookie SSR, inline head script & \textbf{Inline \lstinline|<ThemeInitScript />|} injected before hydration in \lstinline|CommonHead| (Pages Router) and in \lstinline|CommonScripts| (App Router) & \lstinline|packages/mui/src/components/| \lstinline|ThemeContextProvider/ThemeInitScript.tsx| \\
        SSR / hydration & --- & \lstinline|data-theme| set synchronously in \lstinline|<head>| before React mounts, so first paint already matches user preference & \lstinline|CommonHead.tsx|, \lstinline|CommonScripts.tsx| \\
        \bottomrule
    \end{tabular}
\end{table}

\subsubsection{The critical piece --- sync MUI theme with global CSS variables}

The codebase mixes \textbf{MUI \lstinline|sx| / \lstinline|styled|} and \textbf{plain CSS modules}. To keep them in sync we run a single source of truth (CSS variables) and let both worlds resolve to it.

\textbf{Algorithm}

\begin{enumerate}
    \item \lstinline|global-css| declares two palettes as CSS custom properties, scoped by \lstinline|:root[data-theme]|.
    \item \lstinline|global-css/colorPalette.ts| re-exports JS tokens whose value is \lstinline|'var(--red-800)'|, \lstinline|'var(--gray-100)'|, etc. --- \textbf{not} hex.
    \item \lstinline|packages/mui/src/theme.tsx| builds the MUI theme with \lstinline|cssVariables: true| (new in MUI v7) and feeds \lstinline|palette.mode| from the runtime.
    \item \lstinline|ThemeContextProvider| toggles \lstinline|document.documentElement.dataset.theme| --- this flips the CSS variables globally, which in turn re-tints every MUI component and every CSS module simultaneously.
\end{enumerate}

\begin{lstlisting}[caption={packages/global-css/src/styles/palettes/custom-light-variables.css}]
:root,
:root[data-theme='light'],
.light-mode {
  --palette-mode: 'light';
  --inverse-mode: 'dark';

  --rakuten-logo: var(--red-800);
  --merchant-center-logo: '#4D4D4D';
}
\end{lstlisting}

\begin{lstlisting}[caption={packages/global-css/src/colorPalette.ts}]
export const colorPalette = {
  mode: 'var(--mode)',
  reverseMode: 'var(--inverse-mode)',

  white: 'var(--white)',
  gray100: 'var(--gray-100)',
  // ...
  red800: 'var(--red-800)',
  rakutenLogoColor: 'var(--rakuten-logo)',
  merchantCenterLogo: 'var(--merchant-center-logo)',
};
\end{lstlisting}

\subsubsection{Fighting the Flash Of Unstyled Content}

Without pre-hydration work, the browser paints the default (light) palette first, then React swaps to dark 200--800\,ms later. Solution: a \textbf{synchronous inline script} that sets \lstinline|data-theme| before the CSS parses.

\begin{lstlisting}[caption={packages/mui/src/components/ThemeContextProvider/ThemeInitScript.tsx}]
'use client';

import React from 'react';

export function ThemeInitScript() {
  return (
    <script
      id="theme-init"
      dangerouslySetInnerHTML={{
        __html: `
          document.documentElement.dataset.theme =
            localStorage.getItem("theme") ||
            (window.matchMedia("(prefers-color-scheme: dark)").matches ? "dark" : "light");
        `,
      }}
    />
  );
}
\end{lstlisting}

\begin{itemize}
    \item \textbf{Pages Router}: injected in \lstinline|CommonHead| (which is used by both \lstinline|_document| and \lstinline|_documentEmotionCache|).
    $\to$ commit \lstinline|1bde5fcfed| \emph{``fix(theme-init): Injected ThemeInitScript in CommonHead rather than \_document so it gets injected in both \_document \& \_documentEmotionCache''}.
    \item \textbf{App Router}: injected in \lstinline|CommonScripts|.
\end{itemize}

\subsubsection{POC --- MUI v5 vs MUI v7}

A small POC (\lstinline|feature/poc-mui-upgrade-agent-os|) was run first to confirm:
\begin{itemize}
    \item v5 has \textbf{no} first-class \lstinline|cssVariables| flag $\to$ JS $\leftrightarrow$ CSS sync would require a hand-rolled \lstinline|CssVarsProvider| wrapper.
    \item v7 exposes \lstinline|cssVariables: true| on \lstinline|createTheme|, plus a much cleaner \lstinline|variants| API to replace \lstinline|ownerState| conditionals.
    \item v5 $\to$ v7 is a prerequisite, so the migration was folded into the same delivery.
\end{itemize}

\subsection{Phase 02 --- Technical Implementation}

\subsubsection{Upgrade to MUI v7}

\textbf{Version bump} --- every \lstinline|@mui/*| package pinned to the v7 line:

\begin{lstlisting}[]
"@mui/material":         "^7.3.8",
"@mui/icons-material":   "^7.3.8",
"@mui/system":           "^7.3.8",
"@mui/lab":              "^7.0.0",
"@mui/material-nextjs":  "^7.3.8",
"@mui/x-date-pickers":   "^7.29.4"
\end{lstlisting}

\textbf{Import rules updated across the codebase} (commit \lstinline|6903d9e359| \emph{``refactor(mui): Updated theme related imports in MUI components to follow MUI V7 API''}):

\begin{itemize}
    \item Everything theme-related --- \lstinline|useTheme|, \lstinline|createTheme|, \lstinline|styled|, \lstinline|alpha|, \lstinline|Theme|, \lstinline|PaletteMode| --- moves to \lstinline|@mui/material/styles|.
    \item Runtime components/props (\lstinline|SxProps|, \lstinline|useMediaQuery|, \lstinline|Drawer|, etc.) stay on \lstinline|@mui/material| (top-level export only, deep imports are banned).
\end{itemize}

\textbf{Breaking changes handled}

\begin{itemize}
    \item \lstinline|Grid| $\to$ \lstinline|GridLegacy| (9 files: \lstinline|TextStyles.stories|, \lstinline|ToggleButtonGroup.stories|, \lstinline|Title.stories|, \lstinline|StyledModal.stories|, \lstinline|SingleForm|, \lstinline|Dropzone|, \lstinline|Dropzone.style|, \lstinline|Radio.stories|, \lstinline|LoadingScreen|).
    \item \lstinline|ownerState| variant conditionals rewritten into the new \lstinline|variants: [{ props, style }]| array (see \lstinline|Button|, \lstinline|Radio|, \lstinline|Checkbox| in \lstinline|theme.tsx|).
    \item \lstinline|@mui/styles| removed (deprecated in v6, gone in v7) --- from both \lstinline|@rakutenfrance/mui| and \lstinline|@rakutenfrance/nextjs-app-router|.
    \item Deep import \lstinline|Box/Box| fixed in \lstinline|ToggleButton.stories|.
    \item \lstinline|react-is| pinned to \lstinline|^18.3.1| in root \lstinline|pnpm.overrides| to keep MUI v7 happy on React 18 consumers.
    \item \lstinline|x-date-pickers|: \lstinline|calendarPickerClasses|/\lstinline|calendarOrClockPickerClasses| removed, only \lstinline|pickersPopperClasses| kept.
\end{itemize}

Full migration diary: \lstinline|MUI_V7_MIGRATION_SUMMARY.md| at the repo root and \lstinline|agent-os/specs/2026-02-16-1530-mui-v5-to-v7-upgrade/|.

\subsubsection{Theme tokens --- \texorpdfstring{\lstinline|global-css|}{global-css} becomes the source of truth}

\textbf{New files} (commit \lstinline|a03deb6b7b| \emph{``feat(global-css): Added Dark and Light Palettes' Colors''}):

\begin{lstlisting}[]
packages/global-css/src/styles/palettes/
|-- custom-dark-variables.css      # Rakuten-specific, dark
|-- custom-light-variables.css     # Rakuten-specific, light
|-- dark-palette.css               # Full gray/red/orange/... ramps, dark
`-- light-palette.css              # Full gray/red/orange/... ramps, light
\end{lstlisting}

\begin{itemize}
    \item \lstinline|light-palette.css| binds all \lstinline|--gray-*|, \lstinline|--red-*|, \lstinline|--orange-*|, \lstinline|--overlay-*|, plus club status colors, under \lstinline|:root, :root[data-theme='light'], .light-mode|.
    \item \lstinline|dark-palette.css| re-binds the \textbf{same} variable names to inverted values under \lstinline|:root[data-theme='dark'], .dark-mode|. Nothing else has to know which mode is active.
    \item \lstinline|custom-*-variables.css| layer app-level tokens (\lstinline|--rakuten-logo|, \lstinline|--merchant-center-logo|, \lstinline|--palette-mode|, \lstinline|--inverse-mode|).
    \item Both palettes are pulled in via \lstinline|styles/colors.css|, which is imported by \lstinline|index.global.css| and consumed by every app.
\end{itemize}

\textbf{\texorpdfstring{\lstinline|colorPalette.ts|}{colorPalette.ts} --- JS mirror of the CSS vars}

Previously lived in \lstinline|packages/mui/src/colorPalette.ts| with hex values. Now:

\begin{itemize}
    \item Moved into \lstinline|packages/global-css/src/colorPalette.ts|.
    \item Every value replaced with a \lstinline|var(--*)| reference --- see commit \lstinline|161dc47ba2| \emph{``feat(mui-theme): Added CSS Variables support to the Theme Color Palette''}.
    \item \lstinline|global-css| now builds \& publishes as a real package (\lstinline|build:esm|, \lstinline|build:cjs|, \lstinline|exports|, \lstinline|main|, \lstinline|types|) instead of shipping raw source only.
\end{itemize}

Result: \textbf{every downstream \lstinline|.style.ts| and \lstinline|.module.css| file automatically gains dark-mode support} as soon as they import from \lstinline|@rakutenfrance/global-css|.

\subsubsection{The \texorpdfstring{\lstinline|ThemeContextProvider|}{ThemeContextProvider}}

Single React Context wired into the MUI provider tree (commit \lstinline|3407a9770b| \emph{``Implemented Light/Dark Mode Support for all next-common''} + fix \lstinline|e147d96c23| \emph{``fix: persist dark mode on re-renders by syncing data-theme''}).

\begin{lstlisting}[caption={packages/mui/src/components/ThemeContextProvider/ThemeContextProvider.tsx}]
'use client';

import { CssBaseline, useMediaQuery } from '@mui/material';
import { ThemeProvider, type PaletteMode } from '@mui/material/styles';
import React, { useEffect, useMemo, useState } from 'react';
import { globalTheme } from '../../theme';
import { AppThemeContext } from './themeContext';

export const ThemeContextProvider: React.FC<{
  children: React.ReactNode;
  mode?: PaletteMode;
}> = ({ children, mode }) => {
  const prefersDark = useMediaQuery('(prefers-color-scheme: dark)');

  const [manualMode, setManualMode] = useState<PaletteMode | null>(null);

  useEffect(() => {
    const stored =
      (localStorage.getItem('theme') as PaletteMode | null) ||
      (document.documentElement.dataset.theme as PaletteMode | null);
    if (stored) setManualMode(stored);
  }, []);

  const currentMode: PaletteMode = mode ?? manualMode ?? (prefersDark ? 'dark' : 'light');

  const toggleMode = () => {
    const newMode = currentMode === 'light' ? 'dark' : 'light';
    localStorage.setItem('theme', newMode);
    setManualMode(newMode);
  };

  const theme = useMemo(
    () => globalTheme(undefined, undefined, undefined, currentMode),
    [currentMode]
  );

  useEffect(() => {
    document.documentElement.dataset.theme = currentMode;
  }, [currentMode]);

  const appThemeContext = useMemo(
    () => ({ mode: currentMode, toggleMode }),
    [currentMode]
  );

  return (
    <AppThemeContext.Provider value={appThemeContext}>
      <ThemeProvider theme={theme}>
        <CssBaseline />
        {children}
      </ThemeProvider>
    </AppThemeContext.Provider>
  );
};
\end{lstlisting}

\begin{itemize}
    \item Reads persisted mode after mount (no SSR mismatch --- the DOM has already been set by \lstinline|ThemeInitScript|).
    \item Exposes \lstinline|useThemeContext()| $\to$ \lstinline|{ mode, toggleMode }| for any component that needs to render a theme toggle.
    \item Feeds \lstinline|palette.mode| into \lstinline|globalTheme(color, htmlFontSize, pxToRem, currentMode)|.
\end{itemize}

\lstinline|MUIProvider| and \lstinline|AppRouterMUIProvider| now delegate to \lstinline|ThemeContextProvider| instead of wrapping their own \lstinline|ThemeProvider|, so every app that already uses \lstinline|MUIProvider| gets dark mode for free.

\subsubsection{\texorpdfstring{\lstinline|globalTheme|}{globalTheme} --- CSS variables enabled}

\begin{lstlisting}[caption={packages/mui/src/theme.tsx}]
const theme: CustomThemeOptions = {
  pxToRemUtils: pxToRem,
  cssVariables: true,
  typography: {
    // ...
  },
  palette: {
    mode: paletteMode as PaletteMode,
    // ...
  },
  // ...
};
\end{lstlisting}

\begin{itemize}
    \item New \lstinline|_paletteMode| parameter $\to$ drives \lstinline|palette.mode|.
    \item \lstinline|palette.mode| now passed through so MUI generates variable-based classnames (\lstinline|--mui-palette-*|).
    \item All hardcoded palette entries replaced with \lstinline|colorPalette.*| (which are \lstinline|var(--*)|).
    \item Text/background/action/divider filled in so v7 components (\lstinline|Alert|, \lstinline|Chip|, \lstinline|Snackbar|, \lstinline|LinearProgress|) resolve correctly in both modes.
    \item \lstinline|Button|, \lstinline|Radio|, \lstinline|Checkbox|, etc. --- \lstinline|ownerState| conditionals migrated to the v7 \lstinline|variants| array while doing the palette rewrite.
\end{itemize}

\textbf{Known caveat} (documented in commit \lstinline|3407a9770b|): MUI's \lstinline|alpha()| helper does not work on CSS variables, so any component that needed transparency was converted to use \lstinline|--overlay-*| tokens instead.

\subsubsection{Storybook}

Two decorators + a global toolbar toggle (commit \lstinline|3407a9770b| + \lstinline|e21385b506| \emph{``fix(storybook): Added CacheProvider to Preview''}).

\begin{lstlisting}[caption={.storybook/decorators/story-theme-decorators.tsx}]
import type { Decorator } from '@storybook/react';
import { PaletteMode } from '@mui/material';
import { ThemeContextProvider } from '@rakutenfrance/mui';

export const withCssVarsTheme: Decorator = (Story, context) => {
  const mode = context.globals.theme ?? 'dark';
  document.documentElement.setAttribute('data-theme', mode);

  return <Story {...context.args} />;
};

export const withMuiTheme: Decorator = (Story, context) => {
  const mode = (context.globals.theme ?? 'dark') as PaletteMode;

  document.documentElement.setAttribute('data-theme', mode);

  return (
    <ThemeContextProvider mode={mode}>
      <Story {...context.args} />
    </ThemeContextProvider>
  );
};
\end{lstlisting}

\begin{itemize}
    \item \lstinline|preview.tsx| registers both decorators, adds a \lstinline|theme| globalType with a sun/moon toolbar, and now imports the raw \lstinline|@rakutenfrance/global-css/src/index.global.css| so the palettes are available in the iframe.
    \item \lstinline|preview.css| no longer forces \lstinline|background: #fff|; it uses \lstinline|var(--white)| so it inherits the active palette.
    \item Added \lstinline|CacheProvider| (Emotion) around every story so SSR-style class names stay stable across the toggle.
\end{itemize}

\begin{figure}[h]
    \centering
    \fbox{\parbox{0.8\textwidth}{\centering \vspace{2cm} Placeholder screenshot \vspace{2cm}}}
    \caption{Storybook toolbar showing Light/Dark toggle and a component rendering in both modes.}
    \label{fig:storybook-toggle}
\end{figure}

\subsubsection{Shared SVGs --- circular dependency \& Rakuten logo}

Before this branch, \lstinline|@rakutenfrance/mui| exported \lstinline|colorPalette|, and \lstinline|@rakutenfrance/shared-svg| imported it. But \lstinline|mui| also depends on \lstinline|shared-svg| for icons --- so a proper build of the palette would have created a \textbf{circular dependency}.

\textbf{Fix} (commit \lstinline|53d853cfd2| \emph{``refactor: resolved circular dependency between @rakutenfrance/mui \& @rakutenfrance/shared-svg packages''}):

\begin{itemize}
    \item \lstinline|colorPalette| extracted to \lstinline|@rakutenfrance/global-css|.
    \item \lstinline|@rakutenfrance/shared-svg| now depends on \lstinline|@rakutenfrance/global-css| (leaf package, no back-edge).
    \item All SVGs default their fill to a token: \lstinline|fill = colorPalette.gray800| (icons) or \lstinline|fill = colorPalette.rakutenLogoColor| (logos).
    \item \lstinline|SvgImg| / \lstinline|SvgIcon| base components accept and forward \lstinline|fill|.
\end{itemize}

\textbf{Rakuten \& ClubR logos} (commit \lstinline|37c6a257c2| \emph{``fix(logos): Improve Rakuten and ClubR Logos for Dark/Light Mode''}):

\begin{itemize}
    \item \lstinline|RakutenIchiba.tsx| refactored to accept \lstinline|fill| (defaulting to \lstinline|colorPalette.rakutenLogoColor| = \lstinline|var(--rakuten-logo)|).
    \item In light mode \lstinline|--rakuten-logo| $\to$ \lstinline|var(--red-800)| (\lstinline|#BF0000|); in dark mode it flips to \lstinline|var(--color-0)| (white).
    \item Duplicate logo removed from \lstinline|header|; single source in \lstinline|shared-svg|.
\end{itemize}

\textbf{Merchant Center logos} (commit \lstinline|43c9acc252| \emph{``enhance(merchantCenterHeader): Added Dark/Light Mode Support to Merchant Center Header''}):

\begin{itemize}
    \item \lstinline|MerchantCenterLogo|, \lstinline|MerchantCenterLogoMobile| converted \textbf{from raster URLs to inline SVG React components} so they can respect \lstinline|--merchant-center-logo|.
\end{itemize}

\begin{figure}[h]
    \centering
    \fbox{\parbox{0.8\textwidth}{\centering \vspace{2cm} Placeholder screenshot \vspace{2cm}}}
    \caption{Rakuten Ichiba + Merchant Center logos in both modes.}
    \label{fig:logos-both-modes}
\end{figure}

\subsubsection{Dark mode in Header \& Footer}

\textbf{Header} --- see commits \lstinline|f0a243682f|, \lstinline|37c6a257c2|, \lstinline|622bf98ac4|, \lstinline|223d1367cb|:

\begin{itemize}
    \item Every \lstinline|.style.ts| file that referenced \lstinline|colorPalette| from \lstinline|@rakutenfrance/mui| now pulls it from \lstinline|@rakutenfrance/global-css|.
    \item Hardcoded hex in \lstinline|Highlights|, \lstinline|PreHeader Renderer|, \lstinline|Nav|, \lstinline|Menu|, \lstinline|SearchCtn|, \lstinline|PriceClub|, tooltips, etc. $\to$ replaced with CSS variables.
    \item \lstinline|PreHeader| inline HTML string builder: \lstinline|#BF0000| $\to$ \lstinline|var(--red-800)|, \lstinline|#FFFFFF| $\to$ \lstinline|var(--white)|, chevron SVGs default to \lstinline|colorPalette.white|.
    \item Menu categories icon color was intentionally frozen (fix \lstinline|223d1367cb|) to keep the \texteuro-sign color on Super Bons Plans consistent regardless of mode.
    \item \lstinline|preHeader| icons re-styled to work on both backgrounds.
\end{itemize}

\begin{lstlisting}[caption={packages/header/src/PreHeader/src/Renderer/longStringValues.ts}]
'padding:0 16px;' +
'background-color: var(--red-800);' +
'color:var(--white);' +
\end{lstlisting}

\textbf{Footer}:

\begin{itemize}
    \item Every \lstinline|.module.css| cleaned of trailing hardcoded colors.
    \item \lstinline|RakutenGroupLinksLogo| swapped its \lstinline|<img src=...>| for \lstinline|<RakutenIchiba />| so the corporate logo inherits the theme.
\end{itemize}

\begin{figure}[h]
    \centering
    \fbox{\parbox{0.8\textwidth}{\centering \vspace{2cm} Placeholder screenshot \vspace{2cm}}}
    \caption{Header (with preheader + highlights + search) and Footer, both modes.}
    \label{fig:header-footer-both-modes}
\end{figure}

\subsubsection{Dark mode in the rest of the shared components}

Commits \lstinline|3407a9770b|, \lstinline|1864a50d8d|, \lstinline|a0d9b2febc|, \lstinline|5af5869de2|, \lstinline|1431b687bc|, \lstinline|fe9cf4e8b4|:

\begin{longtable}{@{}p{4.3cm}p{9.7cm}@{}}
    \caption{Dark-mode updates applied across shared component packages} \label{tab:shared-components} \\
    \toprule
    \textbf{Package} & \textbf{What changed} \\
    \midrule
    \endfirsthead
    \toprule
    \textbf{Package} & \textbf{What changed} \\
    \midrule
    \endhead
    \bottomrule
    \endfoot
    \lstinline|mui/CustomDrawer|, \lstinline|mui/Drawer|, \lstinline|mui/DrawerLoiLemaire|, \lstinline|mui/SellerDrawer| & v7 import fix (\lstinline|Theme, styled| from \lstinline|@mui/material/styles|); backgrounds $\to$ \lstinline|colorPalette.gray100|; borders $\to$ \lstinline|colorPalette.gray200|; icons $\to$ \lstinline|colorPalette.gray700| \\
    \lstinline|mui/Dialog|, \lstinline|mui/DialogAlert|, \lstinline|PreHeaderDialog| & v7 imports, palette tokens \\
    \lstinline|mui/ImageGallery|, \lstinline|product-image/DesktopImageGallery| & Background $\to$ \lstinline|colorPalette.white| so it inverts in dark mode \\
    \lstinline|mui/Stepper| & Step icon colors driven by palette variables \\
    \lstinline|map| package (\lstinline|RakutenMap|, \lstinline|MultiplePickUpMarker|, \lstinline|PickUpSelectionPanel|, \lstinline|SearchInZone|, \lstinline|CloseButton|, etc.) & Every direct hex replaced; \lstinline|colorPalette| moved to \lstinline|@rakutenfrance/global-css| import; markers respect palette \\
    \lstinline|eshop-header| & New responsive \lstinline|sx| (no more \lstinline|theme.breakpoints.down| closures) + palette tokens \\
    \lstinline|merchant-center-header| & Logo swap + palette tokens for tabs/notifications/monetization \\
    \lstinline|seller-drawer|, \lstinline|strapi|, \lstinline|product-image|, \lstinline|product-description|, \lstinline|paymentMeanDetailsModalV5|, \lstinline|tfa-settings|, \lstinline|dashboard|, \lstinline|addToCartMui|, \lstinline|advertCrewPartnerCard|, \lstinline|sponsor-ads|, \lstinline|crew-partner-modal|, \lstinline|couponClipboard|, \lstinline|crossedPrice|, \lstinline|card*|, \lstinline|carousel|, \lstinline|button|, \lstinline|alert|, \lstinline|avatarLogo|, \lstinline|savebar|, \lstinline|passwordField|, \lstinline|progress|, \lstinline|radioContainer|, \lstinline|rating|, \lstinline|switch|, \lstinline|stepper|, \lstinline|snackbar|, \lstinline|sellerLabel|, \lstinline|productState|, \lstinline|optionsMenu|, \lstinline|link|, \lstinline|label*|, \lstinline|imageGallery|, \lstinline|hyperlink|, \lstinline|helpLinkList|, \lstinline|groupedDropzone|, \lstinline|filter|, \lstinline|feedbackButton|, \lstinline|fab|, \lstinline|extensionClub|, \lstinline|dropzone|, \lstinline|datesSelector|, \lstinline|customDrawerMobile|, \lstinline|customBaseTable|, \lstinline|countryFlag|, \lstinline|copyTooltip|, \lstinline|clubRankProgress|, \lstinline|clubRCagnotte|, \lstinline|clubBalanceCard|, \lstinline|clubBalanceBar|, \lstinline|boxLink|, \lstinline|blockTypeB|, \lstinline|breadcrumb|, \lstinline|advantage|, \lstinline|actionBox| & Same recipe: import \lstinline|colorPalette| from \lstinline|@rakutenfrance/global-css|, remove forced \lstinline|#fff| / \lstinline|#000|, use \lstinline|colorPalette.white| / \lstinline|colorPalette.gray1000| \\
\end{longtable}

Forced background/foreground purge --- commit \lstinline|1431b687bc| \emph{``refactor(style): set backgroundColor to colorPalette.white''} and \lstinline|fe9cf4e8b4| \emph{``style: Updated forced white and black colors to dynamic css variables values''}.

\begin{figure}[h]
    \centering
    \fbox{\parbox{0.8\textwidth}{\centering \vspace{2cm} Placeholder screenshot \vspace{2cm}}}
    \caption{Drawer + Dialog + Map + ProductCard in both modes.}
    \label{fig:components-both-modes}
\end{figure}

\subsubsection{Build \& release to GitHub Packages}

\begin{itemize}
    \item \lstinline|@rakutenfrance/global-css| now publishes real ESM + CJS artifacts (\lstinline|dist/esm|, \lstinline|dist/cjs|, \lstinline|main|, \lstinline|types|, \lstinline|exports| entries) so downstream apps can \lstinline|import { colorPalette } from '@rakutenfrance/global-css'| without pulling raw source.
    \item \lstinline|@rakutenfrance/nextjs|: \lstinline|@mui/styles| peer dep dropped, \lstinline|@mui/material-nextjs| upgraded to \lstinline|^7.3.8|.
    \item \lstinline|@rakutenfrance/nextjs-app-router|: same --- \lstinline|@mui/styles| removed.
    \item \lstinline|apps/nextjs| (POC app) had Turbopack configured with a set of stubs (\lstinline|fs-stub|, \lstinline|prom-client-stub|, \lstinline|html-context-stub|, \lstinline|empty-module|) to unblock server-only modules under Turbopack.
    \item Released as pre-release tags: \lstinline|v14.2.0-rfr-4547.0| $\to$ \lstinline|.1| $\to$ \lstinline|.2| $\to$ \lstinline|.3| $\to$ \lstinline|.4|.
\end{itemize}

\subsection{Phase 03 --- Rollout to fo-apps (next steps)}

For each downstream app that already consumes \lstinline|@rakutenfrance/mui| + \lstinline|@rakutenfrance/header| + \lstinline|@rakutenfrance/footer|:

\begin{enumerate}
    \item \textbf{Upgrade the app to MUI v7} --- bump \lstinline|@mui/*| peers, apply the same import rules (\lstinline|Theme|, \lstinline|styled|, \lstinline|useTheme| from \lstinline|@mui/material/styles|), migrate \lstinline|Grid| $\to$ \lstinline|GridLegacy| where needed.
    \item \textbf{Bump the \lstinline|@rakutenfrance/*| packages} to the \lstinline|14.2.0-rfr-4547.x| line.
    \item \textbf{Fix imports} --- anywhere the app imports \lstinline|colorPalette| from \lstinline|@rakutenfrance/mui|, redirect it to \lstinline|@rakutenfrance/global-css|.
    \item \textbf{Wrap the app} in \lstinline|MUIProvider| (Pages Router) or \lstinline|CommonProviders| / \lstinline|AppRouterMUIProvider| (App Router) --- \lstinline|ThemeContextProvider| is already inside.
    \item \textbf{Inject \lstinline|ThemeInitScript|} --- automatic if the app uses \lstinline|CommonHead| (Pages) or \lstinline|CommonScripts| (App Router). Otherwise add it explicitly to the app's \lstinline|<head>|.
    \item \textbf{Add the theme toggle UI} --- call \lstinline|useThemeContext()| from \lstinline|@rakutenfrance/mui| in the header/user-menu component to expose \lstinline|toggleMode|.
    \item \textbf{Purge app-level hardcoded hex} in any local \lstinline|.style.ts| / \lstinline|.module.css| --- replace with \lstinline|colorPalette.*| or \lstinline|var(--*)|.
    \item \textbf{QA on the dev environment} --- visual regression on Storybook + manual smoke on the app's main flows in both modes.
\end{enumerate}

\subsection{Follow-ups \& known caveats}

\begin{itemize}
    \item MUI's \lstinline|alpha()| helper is incompatible with CSS variables --- any transparency has to be pre-baked in \lstinline|--overlay-*| tokens.
    \item Merchant Center menu icons intentionally frozen to keep Super Bons Plans euro sign colored (fix \lstinline|223d1367cb|).
    \item Storybook currently defaults to \lstinline|light|; can be revisited to auto-follow the OS setting.
    \item The v5 $\to$ v7 migration diary --- including third-party impact (\lstinline|@tolgee/web|, \lstinline|@mui/x-date-pickers|) --- is captured in:
    \begin{itemize}
        \item \lstinline|MUI_V7_MIGRATION_SUMMARY.md|
        \item \lstinline|agent-os/specs/2026-02-16-1530-mui-v5-to-v7-upgrade/{plan,shape,standards,references,breaking-changes,test-plan}.md|
    \end{itemize}
\end{itemize}

\clearpage
